\subsection{TOPOLOGICAL DIRECT SUM OF STABLE SUBGROUPS}

Since the study of commutative groups with operators is equivalent to the study of modules we shall sometimes allow ourselves to use the terminology proper to modules for arbitrary commutative groups with operators; thus we may speak of linear mappings instead of homomorphisms of commutative groups with operators, and we may use the word projector to denote an idempotent endomorphism of a commutative group with operators.

If a commutative topological group \(\mathbf{E}\) with operators (written additively) is the direct sum of a finite family \((\mathbf{M}_i)_{1 \leqslant i \leqslant n}\) of stable subgroups, then the canonical bijection \((x_i) \to \sum_{i=1}^{n} x_i\) of the product group \(\prod_{i=1}^{n} \mathbf{M}_i\) onto \(\mathbf{E}\) is continuous, but is not necessarily a homeomorphism.

DEFINITION 1. Let E be a commutative topological group with operators, and let \((\mathbf{M}_i)_{1 \leqslant i \leqslant n}\) be a finite family of stable subgroups of E, such that E is the direct sum of the \(\mathbf{M}_i\) . Then E is said to be the topological direct sum of the \(\mathbf{M}_i\) if the canonical mapping \((x_i) \to \sum_{i=1}^{n} x_i\) of the product group \(\prod_{i=1}^{n} \mathbf{M}_i\) onto E is a homeomorphism (and therefore an isomorphism of topological groups with operators).

PROPOSITION 2. Let E be a commutative topological group with operators which is the direct sum of stable subgroups \(\mathbf{M}_i\) ( \(1 \leqslant i \leqslant n\) ). Let \((p_i)_{1 \leqslant i \leqslant n}\) be the family of projectors associated with the decomposition \(\mathbf{E} = \sum_{i=1}^{n} \mathbf{M}_i\) . Then E is the topological direct sum of the \(\mathbf{M}_i\) if and only if the \(p_i\) are continuous.

For the mapping \(x \to (p_i(x))\) is the inverse of the mapping \((x_i) \to \sum_{i=1}^{n} x_i\) .

Since \(\mathbf{I} = \sum_{i=1}^{n} p_i\) (where \(\mathbf{I}\) denotes the identity mapping of \(\mathbf{E}\) ) it is sufficient for \(n - \mathbf{I}\) of the projectors \(p_i\) to be continuous in order that the \(n\) th should also be continuous.

If E is the topological direct sum of two stable subgroups M, N, then N is said to be a topological complement of M in E; this is the case if and only if the canonical mapping of E/M onto N is an isomorphism of topological groups with operators.

COROLLARY. Let E be a commutative topological group with operators, and let M be a stable subgroup of E. Then the following conditions are equivalent:

\begin{enumerate}
\def\labelenumi{\alph{enumi})}
\item
  M has a topological complement in E.
\item
  There is a continuous projector \(p\) of \(\mathbf{E}\) into \(\mathbf{E}\) such that \(p(\mathbf{E}) = \mathbf{M}\) .
\item
  The identity mapping of M can be extended to a continuous linear mapping of E onto M.
\end{enumerate}

It follows from Proposition 2 that a) implies b), and it is clear that b) implies c). Finally, if p is a continuous linear mapping of E onto M which extends the identity mapping of M, then p is a continuous projector, and the projectors p and 1-p are associated with the direct sum decomposition \(E = M + N\) , where \(N = p^{-1}(\mathbf{o})\) .

Remarks. 1) To avoid confusion, we shall sometimes say that a stable subgroup of E which is a complement of M (in the sense of the structure of group with operators, without a topology) is an algebraic complement of M.

\begin{enumerate}
\def\labelenumi{\arabic{enumi})}
\setcounter{enumi}{1}
\tightlist
\item
  If a Hausdorff commutative topological group with operators is the topological direct sum of a family \((\mathbf{M}_i)_{1\leqslant i\leqslant n}\) of stable subgroups, then each of the subgroups \(\mathbf{M}_i\) is closed in E, for \(\mathbf{M}_i\) is the set of all \(x\in E\) such that \(p_i(x) = x\) (Chapter I, § 8, no. 1, Proposition 2).
\end{enumerate}

PROPOSITION 3. Let E, F be two commutative topological groups with operators, and let u be a continuous linear mapping of E into F. In order that there should exist a continuous linear mapping v of F into E such that \(u \circ v\) is the identity mapping of F (in which case u is said to be right invertible and v is said to be a right inverse of u) it is necessary and sufficient that u is a strict morphism (§ 2, no. 8) of E onto F and that \(u^{-1}(\mathbf{o})\) has a topological complement in E.

The conditions are necessary. For we have then \(u(v(\mathbf{F})) = \mathbf{F}\) and \(a\) fortiori \(u(\mathbf{E}) = \mathbf{F}\) ; furthermore, if \(p = v \circ u\) , then \(p\) is a continuous linear mapping of \(\mathbf{E}\) into itself such that \(p^2 = p\) ; therefore (Corollary to Proposition 2) \(p(\mathbf{E}) = v(u(\mathbf{E})) = v(\mathbf{F})\) has a topological complement \(p^{-1}(\mathbf{o})\) in \(\mathbf{E}\) ; but since \(u(p(x)) = u(x)\) by hypothesis, we have \(u^{-1}(\mathbf{o}) = p^{-1}(\mathbf{o})\) . Finally, the bijective mapping of \(\mathrm{E}/p^{-1}(\mathbf{o})\) onto \(\mathbf{F}\) , associated with \(u\) , is the composition of the bijective mapping of \(\mathrm{E}/u^{-1}(\mathbf{o})\) onto \(v(\mathbf{F})\) , associated with \(p\) , and the restriction of \(u\) to \(v(\mathbf{F})\) ; since \(v\) is continuous, both these mappings are isomorphisms, and therefore \(u\) is a strict morphism of \(\mathbf{E}\) onto \(\mathbf{F}\) .

The conditions are sufficient. For if \(\varphi\) is the canonical homomorphism of E onto \(\mathrm{E}/u^{-1}(\mathbf{o})\) , to say that \(u^{-1}(\mathbf{o})\) has a topological complement M in E is to say that the restriction of \(\varphi\) to M is an isomorphism of M onto \(\mathrm{E}/u^{-1}(\mathbf{o})\) . Since on the other hand \(u=w\circ\varphi\) , where w is an isomorphism of \(\mathrm{E}/u^{-1}(\mathbf{o})\) onto F, we see that the restriction of u to M is an isomorphism of M onto F, and the inverse isomorphism v is therefore such that \(u\circ v\) is the identity mapping of F onto itself.

PROPOSITION 4. Let E, F be two commutative topological groups with operators, and let u be a continuous linear mapping of E into F. In order that there should exist a continuous linear mapping v of F into E such that \(v \circ u\) is the identity mapping of E onto itself (in which case u is said to be left invertible and v is said to be a left inverse of u) it is necessary and sufficient that u is a (topological) isomorphism of E onto u(E), and that u(E) has a topological complement in F.

The conditions are sufficient, for if they are satisfied we obtain a left inverse \(v\) of \(u\) by taking the composition of the isomorphism of \(u(\mathbf{E})\) onto \(\mathbf{E}\) which is the inverse of \(u\) , with a continuous projector of \(\mathbf{F}\) onto \(u(\mathbf{E})\) .

The conditions are necessary. For the relation \(v(u(x)) = x\) shows that \(u^{-1}(\mathbf{o}) = \{0\}\) ; \(u\) is therefore a bijection of \(E\) onto \(u(E)\) , and since the restriction of \(v\) to \(u(E)\) is continuous, it follows that \(u\) is an isomorphism of \(E\) onto \(u(E)\) . On the other hand, if we put \(q = u \circ v\) , then \(q\) is a continuous linear mapping of \(F\) onto \(u(E)\) such that \(q^2 = q\) , which proves (Corollary to Proposition 2) that \(u(E)\) has a topological complement in \(F\) .

